\documentclass[12pt]{ctexart}

\usepackage[a4paper,margin=2.2cm]{geometry}
\usepackage{amsmath,amssymb}
\usepackage{booktabs}
\usepackage{array}
\usepackage{longtable}
\usepackage{xcolor}
\usepackage{hyperref}
\usepackage{enumitem}
\usepackage{setspace}

\setstretch{1.15}
\setlist[itemize]{leftmargin=1.8em}
\setlist[enumerate]{leftmargin=1.8em}
\hypersetup{colorlinks=true,linkcolor=black,urlcolor=blue}

\title{\textbf{RoadMC 技术报告}\\
\large 面向路面病害点云生成、分割训练与下一阶段验证规划}
\author{YQGHL}
\date{2026-06-30}

\begin{document}
\maketitle

\section*{摘要}

RoadMC 是一个面向路面病害识别的点云学习项目，其核心不是单独训练一个分割模型，而是把物理仿真点云生成和病害分割训练做成一个闭环系统。项目先基于路面粗糙度先验、微纹理、病害几何原语和 LiDAR 观测模型合成带标签点云，再用 Swin3D 或 PointMamba 等骨干网络，结合 mHC、Focal/Dice/Edge 损失和 Muon/AdamW 优化策略进行逐点分割训练。当前阶段已经验证了“二分类优先、阈值校准、再迁移到多类别”的路线具备可行性，并得到一组具有代表性的阶段性结果：Swin3D + mHC 路线在当前合成验证条件下，经阈值校准后可达到约 0.7319 的全局点级 mIoU，而 PointMamba + mHC 在短跑实验中达到约 0.5671 的验证 mIoU。

需要特别说明的是，本报告引用的测试结果均属于项目当前阶段的合成小型验证结果，而不是大规模正式 benchmark。历史记录表明，当前多轮快速诊断通常采用小规模生成数据、较小 batch size，以及仅 2 个验证 batch 的快速估计方式；在这种数据量下，模型很容易在极少 batch 上出现过拟合迹象。因此，现阶段数值更适合作为技术路线正确性的阶段性证据，而不应被视为最终泛化结论。下一阶段的重点将放在合成数据扩容、评估口径统一，以及在更大规模数据集上的重新验证。

\section{项目定位与核心问题}

RoadMC 面向的是一个典型但难度很高的工程问题：路面病害点云分割。该问题至少有四个显著特点：

\begin{itemize}
  \item 标注成本高，真实点云难以快速做大规模标注集；
  \item 病害类别多且长尾严重，直接训练 38 类容易塌缩到背景类；
  \item 裂缝、坑槽、剥落等目标通常细长、稀疏、边界模糊；
  \item 点云观测受采样密度、噪声和扫描角度影响明显。
\end{itemize}

因此，项目的关键判断不是“只要换个更强的 backbone 就能解决问题”，而是必须同时控制监督数据的质量和训练路线的稳定性。RoadMC 的核心价值就在于此：它不是把数据生成当成前处理脚本，而是把生成模块本身当作方法主体的一部分，用可控的几何先验、标签分布和观测模型去支撑后续训练。

\section{总体技术路线}

RoadMC 当前包含两条地位同等重要的主线：

\begin{enumerate}
  \item \textbf{点云生成主线}：构造可训练、可控、带逐点标签的路面点云；
  \item \textbf{模型训练主线}：在这些样本上验证病害分割是否能稳定收敛，并逐步扩展到更复杂类别体系。
\end{enumerate}

整个闭环可以概括为：

\begin{center}
\texttt{路面先验 -> 微纹理 -> 病害几何原语 -> LiDAR 观测 -> 合成点云 -> 分割训练 -> 评估校准 -> 多类别迁移}
\end{center}

这一设计的直接好处是：

\begin{itemize}
  \item 可以控制病害几何形态和类别覆盖；
  \item 可以控制点云更像规则网格还是真实扫描；
  \item 可以把“先学前景/背景，再学细分类别”的训练顺序固化下来；
  \item 可以把实验失败定位到数据、损失、优化器、骨干或评估口径，而不是把所有问题混在一起。
\end{itemize}

\section{点云生成模块}

\subsection{设计目标}

点云生成模块的目标不是制作“看起来像点云”的随机数据，而是尽量从“路面形貌 + 病害形变 + 传感器观测”三层结构去构造监督数据，使输出样本在几何上和标签上都服务训练。

\subsection{关键组成}

\begin{itemize}
  \item \textbf{路面基底先验}：利用 ISO 8608 粗糙度谱构造路面宏观高程起伏；
  \item \textbf{微纹理建模}：用分数布朗运动、局部曲率扰动和法向变化增加表面细节；
  \item \textbf{病害几何原语}：在几何层面生成裂缝、坑槽、松散、沉陷、车辙、剥落等病害形态；
  \item \textbf{LiDAR 观测仿真}：加入扫描线重采样、距离衰减、角度抖动和密度控制；
  \item \textbf{标签传播与导出}：把病害区域映射为逐点标签，并导出为可训练的 \texttt{.npz} 样本。
\end{itemize}

\subsection{点云是如何生成的}

从实际代码实现看，RoadMC 的点云生成并不是一步完成，而是一个顺序明确的物理与几何流水线，当前主流程由 \texttt{roadmc/data/synthetic/generator.py} 负责调度，底层数学与力学基元由 \texttt{roadmc/data/synthetic/primitives.py} 实现。

当前一条典型的场景生成链路可以概括为：

\begin{enumerate}
  \item 先确定路面类型，例如沥青或水泥路面；
  \item 用 ISO 8608 PSD 谱合成路面宏观轮廓；
  \item 叠加 fBm 微纹理，增加表面细节；
  \item 按自然顺序选择并施加病害原语，例如车辙、沉陷、坑槽、裂缝等；
  \item 对规则网格点云做 LiDAR 扫描线重采样，模拟非均匀采样密度；
  \item 计算局部曲率、反射强度等几何和观测特征；
  \item 加入 LiDAR 距离噪声、角度抖动、dropout 和边缘混合效应；
  \item 做最近邻标签与特征传递；
  \item 进行体素下采样或定点数重采样；
  \item 做坐标归一化并导出为训练样本。
\end{enumerate}

其中有几个实现细节值得特别指出。

第一，路面宏观形貌来自频域谱合成，而不是简单随机高度扰动。
项目在路面表面生成函数中利用 ISO 8608 粗糙度等级对应的功率谱密度参数，
通过逆傅里叶变换构建路面高度场，再由高度场计算法向量。

第二，病害不是直接在标签图上绘制，而是在几何层面施加变形。比如裂缝使用 Bézier 路径、Voronoi 结构和分形扰动组合，坑槽使用超椭圆凹陷，车辙和沉陷使用高斯形变，波浪拥包使用正弦调制，水泥路面还有破碎板、错台、拱起、露骨等独立原语。

第三，LiDAR 观测不是附属后处理。项目先通过 \texttt{resample\_to\_lidar\_pattern()} 把规则网格重采样为模拟扫描线的非均匀分布，再通过 \texttt{simulate\_lidar\_noise()} 加入球坐标距离噪声、角度抖动、丢点和边缘混合，使点云更接近真实传感器输出。

第四，生成结果不是只有坐标，而是直接输出可训练样本。每个场景最终都会形成 \texttt{points}、\texttt{labels}、\texttt{feats}、\texttt{normals} 和 \texttt{pavement\_type} 等字段，这也是项目后续能直接接入训练与评估模块的原因。

当前项目中，一个场景样本通常包含以下字段：
\texttt{points}、\texttt{labels}、\texttt{feats}、\texttt{normals} 和 \texttt{pavement\_type}。
训练阶段可直接读取这些字段，并在 38 类标签体系和二分类折叠之间切换。

\subsection{当前工程状态}

从项目历史记录看，生成模块已经具备以下基础能力：

\begin{itemize}
  \item 支持合成带逐点标签的路面点云；
  \item 支持 38 类标签体系，并可在训练阶段折叠为二分类；
  \item 支持通过参数控制路面类型、粗糙度等级、点数和输出规模；
  \item 已经增加了面向扩容的脚本入口，例如扩展到约 5000 个场景的目标配置。
\end{itemize}

但当前生成侧仍有一个明确瓶颈：随着场景数和点数升高，生成速度会明显下降。历史记录里已经出现过更大数据集生成中断、只完成部分 \texttt{train} 样本的情况。这意味着后续如果要真正把验证推到更大规模数据上，生成器的断点续跑、分批执行和吞吐优化必须优先解决。

\section{模型训练模块}

\subsection{输入与数据加载}

训练阶段从合成 \texttt{.npz} 中读取归一化 XYZ 坐标、强度、曲率、裂缝边界距离等特征。项目还在数据加载阶段引入 \texttt{valid\_mask} 用于忽略 padding 点，并通过采样逻辑优先保留病害点，尽量降低 batch 中全背景样本过多的问题。

\subsection{骨干网络}

项目当前支持两条主干线：

\begin{itemize}
  \item \textbf{Swin3D}：基于窗口注意力聚合局部上下文，当前是更稳的精度基线；
  \item \textbf{PointMamba-inspired Backbone}：基于 Morton 顺序与扫描式状态更新，当前是更轻量、更快的实验备选主线。
\end{itemize}

需要强调的是，PointMamba 这条线在本项目中已经通过多轮最小训练闭环验证，但它仍属于工程化实验实现，当前的意义更接近“有潜力继续推进”，而不是已经完全替代 Swin3D。

\subsection{mHC、损失与优化器}

项目默认支持并启用 mHC 通道混合模块，用于增强深层特征流动。损失函数方面，当前路线使用：

\begin{itemize}
  \item Focal Loss：处理背景主导与类别不平衡；
  \item Dice Loss：稳定小目标与稀疏病害学习；
  \item Soft Edge Loss：增强裂缝和边界区域的约束。
\end{itemize}

优化器方面，当前默认采用 Muon + AdamW 的混合策略：二维矩阵参数交给 Muon，bias、norm 和一维参数交给 AdamW；当环境中无法使用 Muon 时，可回退到 AdamW。

\subsection{除了 baseline 以外的训练分支}

除了当前最主要的 baseline 分割训练外，项目实际上还保留了两条额外训练分支，
即 \texttt{gan\_enhanced} 和 \texttt{end2end}。
这两条线虽然不是当前阶段的主结果来源，但它们代表了项目更长远的一个方向，
即从“纯合成监督训练”进一步走向“合成到真实风格适配”。

\subsubsection{GAN 分支的作用}

GAN 相关模块位于 \texttt{roadmc/models/gan/} 目录，当前包含：

\begin{itemize}
  \item \textbf{StyleTransferGen}：一个 DGCNN/EdgeConv 风格的生成器，输入为合成点云坐标和法向量，输出坐标残差与法向量调整量；
  \item \textbf{WGANDiscriminator}：一个 PointNet 风格判别器，用于区分输入点云是否更接近目标风格分布。
\end{itemize}

这条路线的出发点是：即使合成点云在几何和标签上已经足够可控，它和真实采集点云之间仍可能存在“风格差异”，例如局部密度分布、法向扰动、边缘细节和噪声特性不完全一致。GAN 分支的目标就是尝试把合成点云再向更真实的风格分布做一次映射。

\subsubsection{gan\_enhanced 模式}

\texttt{gan\_enhanced} 模式采用两阶段思路：

\begin{enumerate}
  \item 先用 WGAN-GP 方式预训练生成器和判别器；
  \item 再冻结或半冻结生成器，用 GAN 风格调整后的点云去训练分割模型。
\end{enumerate}

从 \texttt{train.py} 当前实现看，这一分支已经具备可运行训练逻辑，并且已经与主干的优化器和 backbone 选择能力打通。但它当前更像是“领域适配支线”或“增强实验入口”，而不是当前最稳定、最主要的交付主线。

\subsubsection{end2end 模式}

\texttt{end2end} 模式更进一步，把 GAN 和分割训练交替放在同一训练循环中，形成：

\begin{center}
\texttt{判别器更新 -> 生成器更新 -> 在风格化点云上训练分割模型}
\end{center}

这条路线的意义在于让风格迁移和分割学习形成共适应关系，但它的训练复杂度、稳定性要求和资源开销都会明显高于 baseline。因此，从项目当前阶段来看，这条分支应理解为已经打通的研究入口，而不是当前已经完成系统验证的稳定主线。

\subsubsection{当前对 GAN / end2end 的定位}

基于现有代码和历史记录，当前对这两条分支最合适的表述是：

\begin{itemize}
  \item 它们是项目的重要支线，不应在技术报告里完全缺席；
  \item 它们代表项目向真实域适配和风格迁移扩展的潜在方向；
  \item 但目前项目的核心阶段性结论，仍主要建立在 baseline 分割主线之上；
  \item 在更大规模数据和更统一评估口径准备好之前，GAN / end2end 还不应被当作当前最终主结果来源。
\end{itemize}

\section{为什么先做二分类}

当前项目没有直接把 38 类作为第一优先级，而是采用：

\begin{center}
\texttt{二分类稳定化 -> 阈值校准 -> backbone 迁移 -> 38 类细分}
\end{center}

这条路线的原因很明确：

\begin{itemize}
  \item 38 类任务长尾过强，冷启动时极易塌缩到背景类；
  \item 如果模型连“病害 / 非病害”的边界都学不稳，继续做 38 类只会放大不稳定性；
  \item 二分类阶段更适合先验证损失设计、优化器设置和骨干是否有效；
  \item 二分类得到的稳定 backbone，可为多类别任务提供更合理的初始化。
\end{itemize}

项目历史记录已经证明，二分类加权和阈值校准都是有效的：不加权时模型很容易在短训后全部预测为背景；显式引入病害权重后，验证 mIoU 可以从 0 恢复到非零水平；进一步做 disease 概率阈值扫描后，最终点级指标还能继续提升。

\section{当前结果的来源、口径与限制}

\subsection{结果来源说明}

本报告引用的所有结果，均来自项目当前阶段的\textbf{合成小型验证}，不是外部公开 benchmark，也不是大规模正式数据集上的最终结论。现阶段结果主要来源于两类实验：

\begin{enumerate}
  \item 基于项目内快速诊断脚本的短跑实验；
  \item 基于当前合成数据目录 \texttt{data/synthetic\_output} 的阶段性训练与验证。
\end{enumerate}

其中，当前较稳定的一套正式合成数据配置来自：

\begin{itemize}
  \item 训练集：600 个合成场景；
  \item 验证集：150 个合成场景；
  \item 常见训练设置：\texttt{batch\_size=4}、\texttt{max\_points=512}、\texttt{lr=3e-4}。
\end{itemize}

同时，历史记录中的大量结构诊断和迁移验证，使用的是更小的快速配置，例如：

\begin{itemize}
  \item \texttt{batch\_size=2}
  \item \texttt{max\_points=512}
  \item \texttt{eval\_batches=2}
  \item 训练步数通常为 60、80、100、200 或 300 step
\end{itemize}

\subsection{过拟合说明}

需要特别强调的是：由于当前阶段很多验证都建立在生成的小型测试集和仅 2 个验证 batch 的快速估计上，模型非常容易在极少 batch 上出现过拟合迹象。换句话说，现阶段实验更适合作为“技术路线是否有效”的早期证据，而不是作为“最终泛化能力”的严格结论。

用户提出的担忧是成立的：当前数据规模偏小，验证口径也偏轻量，在一些阶段性测试中，模型确实可能在 2 个 batch 左右就表现出明显的过拟合趋势。这种现象并不意味着技术路线错误，而是说明当前实验环境更像是\textbf{快速验证平台}，还不是最后的正式评测环境。

\subsection{评估口径差异}

项目历史记录还显示，Lightning 日志中的 batch mean mIoU 与全验证集累计混淆矩阵得到的 global mIoU 并不完全一致。因此，现阶段数值必须结合口径一起看，不能把不同评估方式下的分数直接并列比较。

\section{阶段性代表结果}

\subsection{Swin3D 主线}

当前最有代表性的二分类结果来自 Swin3D + mHC 路线。在当前合成验证条件下，基于二分类训练、显式病害权重和阈值校准，可得到如下结果：

\begin{itemize}
  \item 最优阈值约为 $\tau \approx 0.44$；
  \item 全局点级 mIoU 约为 0.7319；
  \item disease IoU 约为 0.7052；
  \item precision 约为 0.9039；
  \item recall 约为 0.7624。
\end{itemize}

这说明项目已经不再停留在“模型是否能学到病害”这个层面，而是进入了“如何把现有可用结果变得更稳定、更可复现”的阶段。

\subsection{PointMamba 主线}

PointMamba + mHC 路线目前的定位不是“已经超过 Swin3D”，而是“明确可训练、速度更快、值得继续推进”。当前一个代表性短跑结果为：

\begin{itemize}
  \item 200 step 二分类短跑；
  \item final val mIoU 约为 0.5671；
  \item avg loss 约为 0.8949。
\end{itemize}

结合更早的结构闭环和长跑观察，这条路线已经说明自身并非无效，但其精度、稳定性和多类别迁移效果仍需要在更大规模数据上继续确认。

\subsection{38 类任务现状}

38 类任务目前仍处于探索阶段。历史记录显示，在当前规模和步数下，38 类从零训练或直接加载二分类预训练后继续短训，都容易出现以下问题：

\begin{itemize}
  \item 退化到背景主导；
  \item 只预测少数错误类别；
  \item 在短跑验证中 mIoU 维持在 0。
\end{itemize}

因此，当前最合理的结论并不是“38 类路线失败”，而是“38 类冷启动太硬，必须建立在更稳的二分类 backbone 和更大的数据规模之上”。

\section{当前项目的阶段性判断}

基于历史记录、现有脚本和结果口径，当前可以比较稳妥地做出以下判断：

\begin{enumerate}
  \item 物理仿真点云生成作为方法主体是有效且必要的；
  \item 二分类优先路线已经被证明是当前最稳妥的训练路线；
  \item Swin3D 是当前更成熟的精度基线；
  \item PointMamba 是有潜力继续推进的轻量化候选路线；
  \item GAN / end2end 分支已经打通训练入口，但目前更适合作为域适配和风格迁移支线，而不是当前主结果来源；
  \item 现阶段主要矛盾不是 backbone 完全不可用，而是数据规模不足、验证过轻和多类别冷启动过硬；
  \item 当前结果均应视为阶段性验证结果，而非最终正式结论。
\end{enumerate}

\section{下一阶段验证计划}

下一阶段的工作重点，不应是继续在当前小型验证条件下反复堆叠试验，而应转向更稳妥的扩容和复验流程。

\subsection{数据扩容}

项目下一步将优先把合成数据规模扩大到约 5000 个场景，并继续完善以下工程能力：

\begin{itemize}
  \item 断点续跑；
  \item 分批生成；
  \item 多进程稳定执行；
  \item 更可控的点数和密度配置。
\end{itemize}

这是后续所有正式验证的前提。

\subsection{重新验证二分类基线}

在扩容后的数据集上，下一步将重新验证并统一记录：

\begin{itemize}
  \item Swin3D + mHC 的二分类结果；
  \item PointMamba + mHC 的二分类结果；
  \item binary threshold 在 0.40--0.50 区间内的小网格扫描；
  \item \texttt{binary\_class\_weights} 的小范围联合搜索；
  \item 全验证集 global mIoU 与 batch mean mIoU 的统一输出。
\end{itemize}

\subsection{重新设计多类别迁移}

在二分类基线稳定之后，再继续推进多类别任务。当前更合理的路线不是直接从零训练 38 类，而是考虑以下中间阶段：

\begin{itemize}
  \item 二分类 -> 4 类；
  \item 二分类 -> 8 类；
  \item 二分类 -> 高频子类集；
  \item 最后再扩展到 38 类。
\end{itemize}

这将比“直接切回 38 类硬训”更符合项目当前的实际状态。

\subsection{继续推进 GAN / 域适配支线}

在 baseline 主线足够稳定之后，GAN 与 end2end 分支可以继续承担两个方向的工作：

\begin{itemize}
  \item 让合成点云在局部几何、法向和噪声分布上更接近真实域；
  \item 作为后续合成到真实迁移的先导实验入口。
\end{itemize}

因此，GAN 分支并不是可以删去的历史遗留代码，而是项目中偏中长期的一条研究扩展线。更合理的推进顺序是先稳住 baseline，再让这条支线承担域适配任务。

\section{结论}

RoadMC 当前最重要的成果，不是已经把 38 类任务完整解决，而是已经明确了一条可持续推进的技术主线：先用物理仿真点云生成构建可控监督，再用二分类优先的策略把模型训稳，最后再逐步迁移到更细粒度的多类别病害分割。

同时也必须明确：本报告中的结果均来自生成的小型测试集验证，且当前阶段存在数据规模偏小、验证 batch 数偏少、容易在极少 batch 上过拟合的问题。因此，这些结果更适合作为当前技术路线正确性的阶段性证据，而不是最终大规模泛化能力的结论。下一步的正式验证将基于扩容后的合成数据集重新进行，以便对模型的稳定性、泛化能力和多类别迁移效果作出更可信的判断。

\end{document}
